\chapter{资料整理与方法流程}

\section{结构设计原则}

项目层实现遵循“最小侵入”原则：公共包继续负责页面样式、目录、图表、参考文献、统一构建，以及 \texttt{thesis-smu-postdoc} 的独立模板身份；博士后研究报告特有的封面、题名页和结尾材料仍放在项目自己的 \texttt{extraTex/} 中维护。这样既能保证单项目分发时的业务独立，又不会把尚未稳定沉淀的个性化正文逻辑提前固化到共享包中。

\section{正文组织方式}

正文采用五章结构：

\begin{enumerate}
  \item 研究背景与问题界定；
  \item 资料整理与方法流程；
  \item 结果展示与示范分析；
  \item 讨论与局限；
  \item 结论与应用建议。
\end{enumerate}

\section{示范流程图}

\begin{figure}[H]
  \centering
  \fbox{%
    \parbox[c][4.2cm][c]{0.78\textwidth}{%
      \centering
      公开示意图占位\\[0.4em]
      数据抽取 $\rightarrow$ 变量整理 $\rightarrow$ 风险分层 $\rightarrow$ 可解释输出
    }%
  }
  \caption{博士后研究报告示例中的风险分层流程示意图}
  \label{fig:postdoc-workflow}
\end{figure}

\section{关键变量分组}

\begin{table}[H]
  \centering
  \caption{公开示例中的变量分组说明}
  \label{tab:variables}
  \begin{tabularx}{0.92\textwidth}{@{}lXX@{}}
    \toprule
    分组 & 核心变量 & 说明 \\
    \midrule
    基线信息 & 年龄、性别、入科方式 & 用于描述研究对象的基础特征 \\
    生理指标 & 血压、血乳酸、呼吸频率 & 用于模拟早期病情评估信息 \\
    干预信息 & 抗感染治疗、液体复苏、机械通气 & 用于说明临床管理流程 \\
    结局指标 & 28 天死亡、ICU 住院时长、器官支持天数 & 用于演示主要结果写法 \\
    \bottomrule
  \end{tabularx}
\end{table}
